\documentclass[tikz,border=6pt]{standalone}
% 공통 프리앰블 — PM Day1 교재 그림 (TikZ standalone)
\usepackage{fontspec}
\setmainfont{Apple SD Gothic Neo}
\setsansfont{Apple SD Gothic Neo}
\setmonofont{Menlo}
\usepackage{amssymb}
\usepackage{tikz}
\usetikzlibrary{arrows.meta,positioning,calc,shapes.geometric,fit,backgrounds,matrix,decorations.pathreplacing}
\definecolor{navy}{HTML}{1F3A5F}
\definecolor{teal}{HTML}{2A7F62}
\definecolor{rust}{HTML}{C8553D}
\definecolor{sand}{HTML}{E8E1D5}
\definecolor{ink}{HTML}{222222}
\definecolor{grey}{HTML}{7A7A7A}
\definecolor{lightgrey}{HTML}{EFEFEF}
\definecolor{navylight}{HTML}{DCE6F2}
\definecolor{teallight}{HTML}{DDF0E8}
\definecolor{rustlight}{HTML}{F6E0DA}
\tikzset{
  every node/.style={font=\small, text=ink},
  box/.style={draw=navy, thick, rounded corners=2pt, align=center, inner sep=5pt, fill=white},
  boxfill/.style={box, fill=navylight},
  boxteal/.style={draw=teal, thick, rounded corners=2pt, align=center, inner sep=5pt, fill=teallight},
  boxrust/.style={draw=rust, thick, rounded corners=2pt, align=center, inner sep=5pt, fill=rustlight},
  boxgrey/.style={draw=grey, thick, rounded corners=2pt, align=center, inner sep=5pt, fill=lightgrey},
  arr/.style={-{Stealth[length=2.5mm]}, thick, navy},
  arrteal/.style={-{Stealth[length=2.5mm]}, thick, teal},
  arrrust/.style={-{Stealth[length=2.5mm]}, thick, rust},
  arrgrey/.style={-{Stealth[length=2.5mm]}, thick, grey},
  lbl/.style={font=\footnotesize, text=grey, align=center},
  src/.style={font=\scriptsize, text=grey, align=left},
  title/.style={font=\normalsize\bfseries, text=navy, align=center},
}

\pgfdeclarelayer{bg1}\pgfdeclarelayer{bg2}\pgfsetlayers{bg1,bg2,main}
\begin{document}
\begin{tikzpicture}[node distance=4mm]
% ---- 좌: 계층 포함 도해 ----
\node[box, minimum width=28mm, minimum height=9mm] (p1) {P1 플랫폼\\{\footnotesize 168억}};
\node[box, minimum width=28mm, minimum height=9mm, right=of p1] (p2) {P2 물류 자동화\\{\footnotesize 194억}};
\node[box, minimum width=28mm, minimum height=9mm, right=of p2] (p3) {P3 변화관리\\{\footnotesize 22억}};
\node[box, minimum width=28mm, minimum height=9mm, below=of p1] (p4) {P4 규제 대응\\{\footnotesize 31억}};
\node[box, minimum width=28mm, minimum height=9mm, right=of p4] (p5) {P5 운영 이관\\{\footnotesize 13억}};
\node[lbl, right=of p5, align=left, text width=28mm] (pnote) {프로젝트 =\\산출물(outputs)을 만든다};
\begin{pgfonlayer}{bg2}
\node[draw=teal, thick, rounded corners=3pt, fill=teallight, fit=(p1)(p3)(p4)(p5)(pnote), inner sep=5mm, inner ysep=7mm, label={[teal, font=\small\bfseries, anchor=north west, xshift=2mm, yshift=-1mm]north west:프로그램 ONE ONDAM — 총 428억, 편익 M+24 연 214억}] (prog) {};
\end{pgfonlayer}
\node[boxgrey, minimum height=9mm, text width=52mm, below=6mm of prog.south west, anchor=north west] (ops) {운영 — 서비스운영팀 (박준영, 서비스데스크 11명)\\{\footnotesize 2027-05-29 이후 연간 운영비 27억 (프로그램 예산 외)}};
\node[lbl, right=6mm of ops, text width=50mm, align=left] {운영 = 반복되고 끝이 없다.\\프로젝트↔운영의 경계는\\착수 시점에 정의한다};
\begin{pgfonlayer}{bg1}
\node[draw=navy, thick, rounded corners=3pt, fill=navylight, fit=(prog)(ops), inner sep=5mm, inner ysep=8mm, label={[navy, font=\small\bfseries, anchor=north west, xshift=2mm, yshift=-1mm]north west:포트폴리오 — 전략 목표 아래 프로젝트·프로그램·운영의 전체 그림 (자원 경쟁의 장)}] (port) {};
\end{pgfonlayer}

% ---- 우: MSP 가치 사슬 ----
\node[boxfill, text width=26mm, right=14mm of port.north east, anchor=north west, yshift=-2mm] (o1) {산출물\\outputs};
\node[boxfill, text width=26mm, below=7mm of o1] (o2) {역량\\capabilities};
\node[boxfill, text width=26mm, below=7mm of o2] (o3) {성과\\outcomes};
\node[boxteal, text width=26mm, below=7mm of o3] (o4) {편익\\benefits};
\draw[arr] (o1) -- (o2); \draw[arr] (o2) -- (o3); \draw[arr] (o3) -- (o4);
\node[lbl, right=3mm of o1, text width=40mm, align=left] {P1: 신규 POS, 주문 허브,\\통합 재고, 가맹 발주 포털};
\node[lbl, right=3mm of o2, text width=40mm, align=left] {실시간 재고 가시성};
\node[lbl, right=3mm of o3, text width=40mm, align=left] {매장·물류의 발주 행동 변화};
\node[lbl, right=3mm of o4, text width=40mm, align=left] {폐기율 절감 43억 (P1+P3)\\물류 처리량 절감 96억 (P1+P2)};
\draw[decorate, decoration={brace, amplitude=3pt}, navy, thick] ($(o1.north west)+(-2mm,0)$) -- ($(o1.south west)+(-2mm,0)$) node[midway, left=2mm, lbl, text=navy, align=right] {PM\\책임};
\draw[decorate, decoration={brace, amplitude=3pt}, teal, thick] ($(o2.north west)+(-2mm,0)$) -- ($(o4.south west)+(-2mm,0)$) node[midway, left=1mm, lbl, text=teal, align=right] {프로그램·\\편익 오너\\책임};
\node[title, above=3mm of o1.north, anchor=south, xshift=18mm] {MSP 가치 사슬};

% 출처
\node[src, below=3mm of port.south west, anchor=north west] {출처: ISO 21502:2020, PMBOK Guide 8판(PMI, 2025), MSP 5판(AXELOS, 2020)의 정의와 사례 문서 『㈜온담F\&B홀딩스』를 바탕으로 재구성.};
\end{tikzpicture}
\end{document}
